\section{Introduction and Goals}

We will build a tool called \texttt{urlget}: a small C program that behaves
like a minimal subset of \texttt{curl}. By the end of this tutorial you will
have a single-file client that can:

\begin{itemize}[leftmargin=1.4em]
    \item Fetch content from \texttt{http://} and \texttt{https://} URLs
    \item Optionally route traffic through Tor via SOCKS5
    \item Follow HTTP redirects (with a configurable limit)
    \item Send custom headers, POST bodies, and basic authentication
    \item Write the response body to a file instead of stdout
    \item Compile clean release builds with debug tracing disabled via
          \texttt{-DNDEBUG}
\end{itemize}

\subsection{Design Principles}

\begin{enumerate}[leftmargin=1.4em]
    \item \textbf{Incremental.} Each step produces a coherent program. You can
          stop after any step and still understand what you have.
    \item \textbf{Readable over clever.} Prefer clear control flow and small
          helpers over dense one-liners.
    \item \textbf{POSIX + OpenSSL.} We use standard sockets for TCP and
          OpenSSL 3 for TLS. No third-party HTTP libraries.
    \item \textbf{Explain the \emph{why}.} Every new block is introduced with
          the problem it solves, not only the syntax.
\end{enumerate}

\subsection{What You Need}

\begin{itemize}[leftmargin=1.4em]
    \item A C compiler (\texttt{gcc} or \texttt{clang})
    \item OpenSSL development headers (for the HTTPS step onward), e.g.\
          \texttt{libssl-dev} on Debian/Ubuntu
    \item Optional: a local Tor daemon listening on \texttt{127.0.0.1:9050}
          if you want to try the SOCKS5 path
\end{itemize}

\subsection{How to Read This Tutorial}

Each development step has the same shape:

\begin{enumerate}[leftmargin=1.4em]
    \item A short motivation: what is missing and why we add it now
    \item Focused snippets of the \textbf{new} code (green background)
    \item The \textbf{full source so far}: unchanged lines in light gray,
          newly added lines in green so they stand out in context
\end{enumerate}

The final step assembles the complete \texttt{urlget.c} matching a real,
usable client.
